% ----------------------------------------------------------
\chapter{Introdução}
% ----------------------------------------------------------


\section{Justificativa}

O agronegócio representa um dos pilares fundamentais da economia brasileira, desempenhando um papel estratégico no desenvolvimento regional. Contudo, a eficiência do setor é frequentemente comprometida pelo elevado índice de desperdício na cadeia de suprimentos de produtos perecíveis. Segundo estudos recentes sobre o fluxo de hortifrúti \cite{ciagro2020}, o Brasil, apesar de ser um dos maiores produtores mundiais, enfrenta gargalos severos podendo somar 50\% de desperdício em toda cadeia logística.

Esse cenário gera um impacto econômico direto no custo do produto. O transporte de vegetais, como o tomate-cereja e o pepino, sem uma triagem prévia adequada, resulta na movimentação de mercadorias fora dos padrões de comercialização. Transportar itens que serão inevitavelmente descartados eleva desnecessariamente os gastos com frete, embalagens e mão de obra, reduzindo a rentabilidade do produtor e a competitividade do setor. Portanto, a automação do processo de classificação na origem torna-se uma solução indispensável para otimizar a cadeia produtiva e garantir que apenas produtos aptos ao consumo sejam transportados.

Apesar da eficácia dos sistemas de visão computacional na triagem de alimentos, as soluções industriais disponíveis no mercado possuem custos elevados, tornando-as inacessíveis para pequenos e médios produtores. O desenvolvimento de uma máquina de baixo custo justifica-se pela necessidade de democratizar o acesso a tecnologias de ponta no campo. Ao focar inicialmente no tomate-cereja, com flexibilidade para expansão a outras culturas como o pepino, este projeto busca oferecer uma alternativa viável e escalável que contribua diretamente para a redução do desperdício e a modernização da agricultura familiar e de pequena escala.
% ----------------------------------------------------------
\section{Objetivos}
% ----------------------------------------------------------

Nesta seção, apresentam-se o objetivo geral e os objetivos específicos que norteiam o desenvolvimento deste trabalho.

\subsection{Objetivo Geral}

Desenvolver e avaliar um protótipo de bancada de inspeção óptica de baixo custo que classifique tomates-cereja quanto ao estádio de maturação, à presença de defeitos e ao tamanho, como prova de conceito para a automação da triagem voltada a pequenos e médios produtores rurais.

\subsection{Objetivos Específicos}

\begin{itemize}
	\item Montar uma bancada de aquisição com câmera fixa, iluminação controlada e fundo escuro, com calibração de branco e calibração métrica por padrão de xadrez, de modo a converter medidas em \textit{pixels} para milímetros;
	\item Implementar a segmentação dos frutos por limiarização de Otsu e medir, para cada tomate, o diâmetro, o comprimento, a largura e a circularidade;
	\item Classificar o estádio de maturação (imaturo, intermediário e maduro) a partir da proporção de cores da superfície do fruto, usando uma tabela de consulta (\textit{Look-Up Table}, LUT) gerada por uma rede neural \textit{Multilayer Perceptron} (MLP) treinada com \textit{pixels} agrupados no espaço CIE L*a*b*;
	\item Construir um \textit{dataset} próprio, coletado na bancada, de tomates sadios e com defeitos visíveis (podridão, rachaduras, danos mecânicos e manchas), e treinar por transferência de aprendizado uma rede neural convolucional ResNet-18 para distinguir frutos sadios de frutos com defeito, com validação cruzada agrupada por tomate;
	\item Integrar as etapas em uma interface homem-máquina (IHM) que exiba o resultado da inspeção em tempo real, indique o destino de cada fruto (descarte ou classe de maturação) e registre as inspeções para posterior análise;
	\item Avaliar o desempenho do sistema em um lote de teste independente, formado por tomates não utilizados no treinamento, por meio de acurácia balanceada, sensibilidade, especificidade e intervalos de confiança, tendo como referência a meta de 90\% de acurácia, além de medir o tempo de processamento por fruto.
\end{itemize}

A separação física dos frutos por atuadores e a integração com uma esteira industrial estão fora do escopo deste trabalho e são indicadas como trabalhos futuros.
